% Abhakseite „Das kann ich“ – Zone nur im Gesamt (\mitzone)
%% TODO Ich-kann-Sätze: Platzhalter wie in den Aufgabendateien
\begin{abhakseite}
\ifdefined\mitzone
\abhakgruppe{Kennst du schon}
\abhak{1}{Stammfunktion bilden und Hauptsatz führen}
\abhak{2}{Nullstellen bestimmen}
\abhak{3}{Schnittstellen zweier Graphen berechnen (gleichsetzen)}
\abhak{4}{Dreiecks-, Trapez- und Kreisformeln}
\abhak{5}{Maßstab und Einheiten}
\abhak{6}{Symmetrie am Term erkennen (gerade und ungerade Anteile)}
\abhak{7}{Stammfunktion bilden und Hauptsatz führen – Fehler finden}
\abhak{8}{Stammfunktion bilden und Hauptsatz führen – selbst rechnen}
\fi
\abhakgruppe{Einheit 1 · Fläche zwischen Graph und x-Achse}
\abhak{9}{Fläche zwischen Graph und x-Achse}
\abhak{10}{Fläche zwischen Graph und x-Achse – Fehler finden, Begründen, Anwendung}
\abhakgruppe{Einheit 2 · Fläche zwischen zwei Graphen}
\abhak{11}{Fläche zwischen zwei Graphen}
\abhak{12}{Fläche zwischen zwei Graphen – Fehler finden, Begründen, Anwendung}
\abhakgruppe{Einheit 3 · Zusammengesetzte Flächen, Maßstab und Volumen}
\abhak{13}{Zusammensetzen, Maßstab, Volumen}
\abhak{14}{Zusammensetzen, Maßstab, Volumen – Fehler finden, Begründen, Anwendung, Darstellungswechsel}
\abhakgruppe{Einheit 4 · Flächenbedingungen}
\abhak{15}{Flächenbedingungen}
\abhak{16}{Flächenbedingungen – Fehler finden, Begründen, Anwendung}
\abhakgruppe{Einheit 5 · Das Integral als Flächenbilanz}
\abhak{17}{Flächenbilanz}
\abhak{18}{Flächenbilanz – Fehler finden, Begründen, Anwendung, Darstellungswechsel}
\end{abhakseite}
